\section{Further exact research questions}
\label{research:section}

The questions below concern identities, finite reductions, and the
structure of the remaining analytic coordinates. They are proposed
problems, not conjectures asserted on numerical evidence. Each specifies
an initial case or a concrete output that could settle the problem.

\begin{question}[Evaluate or relate the two cubic coordinates]
Find a second functional identity that reduces $\Omega$ or $\kappa$
to a prescribed algebra of ordinary zeta derivatives, Gamma or Barnes
values, and elementary constants. The convergent moment in
\eqref{tr3:kappa} and the polylogarithmic integral
\eqref{tr3:Dpolylog} provide two explicit starting points. An equality
between those representations is not itself an ordinary-constant
evaluation. A useful first target is a reflection or multiplication
identity for the whole Dirichlet series $D(A)$, whose second derivative
controls $\kappa$.
\end{question}

\begin{question}[The quartic and higher directional layers]
The cubic theorem uses the degree-two Taylor jet of the holomorphic
remainder $R$. At degree four, compute the degree-three jet subject
to symmetry, the exact $C$ axis, and the Euler relation
\eqref{tr3:stuffleR}. Give an explicit finite set of convergent Gamma
or multiple-Gamma moments spanning all $W^{(4)}_{a,b,c}(0)$.
Determine the exact dimension of the formal jet quotient under the
proved functional relations. This is a finite linear-algebra question
and should be separated from numerical or arithmetic independence
of the resulting constants.
\end{question}

\begin{question}[Systematic cancellation of directional coordinates]
The coefficient of $\kappa$ vanishes on a conic, and cyclic averaging
eliminates it identically. At each higher derivative order, construct
the polynomial map from directions to the residual jet coordinates.
Compute its kernel for permutation sums and for small finite families
of rational directions. The desired output is an exact criterion
for combinations reducible to already known coordinates, with the
smallest number of directions needed within that specified formal
relation space.
\end{question}

\begin{question}[The minimal Gamma-separated independent-order kernel]
Complete Question 8 of the exact-identities report: beyond the entire
interpolant identified here, is there a smallest explicit generating
kernel whose singular part separates through a Gamma factor when
all inner orders move independently? Begin at depth three near
$(s_1,s_2,s_3)=(1,1,1)$ and retain three independent perturbations.
Use the coefficients $C_{k,m}$ in \eqref{eht:C} to compare candidate
factorizations at every intersecting prefix divisor. State the
equivalence relation under which ``smallest'' is meant.
\end{question}

\begin{question}[Nonsymmetric harmonic jets and the coordinate $\eta$]
Equations \eqref{eht:first-eta}--\eqref{eht:second-eta} show precisely
where $\eta(a)$ survives after entire normalization. Determine which
linear combinations of second and third independent order jets cancel
all new ordered coordinates. At rational endpoints, compute a finite
character-projection matrix including conductor logarithms. A first
target is the antisymmetric second jet at depth two with endpoints
$a=1/2$, $b=1$; an ordinary Stieltjes reduction would require an
additional identity beyond the symmetric Gamma generator.
\end{question}

\begin{question}[Reflection and multiplication of normalized primitives]
For the primitive \eqref{eht:primitive} and its nonpositive-inner
extensions, derive exact formulas under $a\mapsto1-a$ and
$a\mapsto qa$. Keep the base endpoint fixed or explicitly transform
it, so additive constants remain determined. Differentiate the whole
spectral identity to obtain log-weighted Gamma and Stieltjes
primitives. Compare the resulting constants with standard Barnes
multiple-Gamma normalizations at the first two derivative orders.
\end{question}

\begin{question}[Minimal depth within the mixed-order output class]
The bound $P+1$ is constructive, not optimal in every example.
Characterize when a mixed input reduces to depth at most $P$, or to
ordinary Hurwitz functions and their order derivatives. Start with
one negative factor and two positive factors at a fixed root-of-unity
alphabet. Any minimality theorem must specify the coefficient field,
the allowed argument transformations, and the set of functional
relations; otherwise a depth count alone is not an obstruction.
\end{question}

\begin{question}[A finite map for character and color cancellations]
Apply root-of-unity projections to the exact compiler before evaluating
periods. Can the positive-depth terms be represented in a finite
quotient whose kernel detects all reductions to the depth-one
Stieltjes formulas \eqref{mixed:DFTStieltjes} at a fixed weight and
derivative order? Repeated colors should be treated within the
confluent formula, so the reduction matrix remains meaningful when
several input colors coincide. An exact row certificate is the
appropriate output.
\end{question}

\begin{question}[Spectral derivatives beyond locally determined constants]
Theorem~\ref{mixed:localtheorem} makes every Laurent constant at a
nonpositive outer integer a local calculation. The next regular
coefficient usually depends on the global Mellin remainder. Classify
linear combinations of mixed kernels for which that remainder cancels
to one or two prescribed orders. The two explicit triple kernels
in Section~\ref{mixed:section} provide a small initial family;
their finite MPL normal forms permit legitimate differentiation
before specialization.
\end{question}

\begin{question}[Simultaneous color and order resonance]
Let several $X_j$ tend to $1$ while $C$ tends to a pole and its
derivatives are taken. Determine a joint completed germ whose local
color logarithms and spectral Laurent coefficients are explicit.
The radial variable must respect the weights,
$X_j\mapsto\rho^{p_j}X_j$. A first target is one negative and two
positive inner orders with two colors approaching $1$ at different
rates. Distinguish a joint analytic continuation from an ordinary
Abel limit of a divergent series.
\end{question}

\begin{question}[Four-factor kernels and higher coincident products]
Use the mixed-order compiler for the integer-index pieces of the
four-factor Fourier zero-frequency constraint. Identify a finite
collection of constrained zeta kernels for the different sign
partitions, then compute the fully coincident fourth digamma finite
part in one fixed endpoint coordinate. The cubic directional formulas
should supply boundary checks. The unresolved global kernel must not
be replaced by an undefined product of singular distributions.
\end{question}

\begin{question}[Sampling sets beyond vanishing logarithmic gaps]
Theorem~\ref{sampling:unique} is sharp against unrestricted geometric
sampling, but a particular set of imaginary exponents may admit much
sparser uniqueness sets. For a prescribed finite frequency set
$\{\Im\alpha_\nu\}$, characterize the sampling sequences on which
its trigonometric leading part is determined. Separate this finite
frequency problem from extensions to asymptotic expansions with
exponentially small terms, where functions invisible to every
power--logarithm coefficient can occur.
\end{question}

\begin{question}[Exact Gaussian certificates]
Resume $S_6$ from the frozen rational target and the precisely specified
relation space, using a strategy that controls sparse fill and records
row provenance. A modular candidate must be lifted and replayed over
the rationals. If the chosen space does not contain the target,
compute a separating functional for that exhausted space and then
identify a mathematically justified new relation family. Treat the
revised $S_8$ candidate separately; numerical agreement alone cannot
promote either to a theorem.
\end{question}

\begin{question}[Formal proof kernels and reusable interfaces]
Formalize the finite root filter, sparse shuffle, repeated-pole Taylor
extraction, geometric/Bernoulli elimination, and polynomial jet maps.
Their analytic interface consists of a reusable Mellin-subtraction
lemma with locally normal parameter dependence and an explicit
reciprocal-Gamma zero. A second interface should formalize prefix
residue cancellation at intersecting divisors. These finite and
analytic components can then support later repository identities
without conflating exact algebra, analytic continuation, and
numerical evaluation.
\end{question}
